\section{Implementation Details}\label{sec:implementation}

To implement the \CPTB{} algorithm, we use the open-source Python library \texttt{dp\_accounting.pld}\footnote{\url{https://github.com/google/differential-privacy/tree/main/python/dp_accounting/pld}}. We extend the class \texttt{dp\_accounting.pld.pld\_mechanism.AdditiveNoisePrivacyLoss} to create a class, \texttt{MixtureGaussianPrivacyLoss} that represents the privacy loss distribution of $\calM_{MoG}$, which can be used along with other tools in the \texttt{dp\_accounting.pld} library to implement \CPTB{}. We discuss our implementation and some challenges here. The \texttt{dp\_accounting.pld} library uses the convention that privacy losses are decreasing; we use the same convention in the discussions in this section for consistency.

\subsection{Extending \texttt{AdditiveNoisePrivacyLoss}}

In order to perform all the necessary computations in \CPTB{}, we need to implement the following methods in \texttt{MixtureGaussianPrivacyLoss}:

\begin{enumerate}
\item A method to compute the CDF of the mixture of Gaussians distribution.
\item A method to compute the privacy loss at $x$.
\item An inverse privacy loss method, i.e. a method which takes $\epsilon$ and computes the smallest $x$ achieving this $\epsilon$.
\end{enumerate}

Given the probabilities and sensitivities $\{p_1, p_2, \ldots, p_k\}$ and $\{c_1, c_2, \ldots, c_k\}$, as well as $\sigma$, the first two can easily be done by just summing the PDFs/CDFs of the Gaussians in the mixture. This takes at most $O(k)$ times the runtime of the corresponding method for the (subsampled) Gaussian mechanism.

The third is more problematic. For the subsampled Gaussian mechanism with sampling probability $p$ and sensitivity $1$, the privacy loss function (under the remove adjacency) is:

\[\ln\left(p \exp\left(\frac{-2x-1}{2\sigma^2}\right) + 1 - p\right).\]

This function is easily invertible. However, if we consider $\calM_{MoG}(\{p, 1-p\}, \{c_1, c_2\})$, the privacy loss at $x$ is:

\[\ln\left(p \exp\left(\frac{-2c_1x-c_1^2}{2\sigma^2}\right) + (1-p)\exp\left(\frac{-2c_2x-c_2^2}{2\sigma^2}\right)\right).\]

Because this function includes the sum of two exponential functions of $x$, it is not easy to invert. We instead use binary search to get the smallest multiple of $\Delta_1$ which achieves the desired privacy loss, where $\Delta_1$ is a parameter we choose that trades off between efficiency and accuracy. That is, if $L$ is the privacy loss function, and we want to compute the inverse privacy loss of $y$, we return $x = \lceil L^{-1}(y) / \Delta_1 \rceil \cdot \Delta_1$.
Note that by overestimating $x$, we also overestimate the privacy loss since we assume the privacy loss is decreasing. Hence this approximation is ``pessimistic,'' i.e. does not cause us to report an $(\epsilon, \delta)$-DP guarantee that is not actually satisfied by $\calM_{MoG}$.

Note that using binary search requires a $O(\log(1/\Delta_1))$ multiplicative dependence on $\Delta_1$, that is not incurred for e.g. the subsampled Gaussian for which we can quickly compute the exact inverse privacy loss. Indeed, we observed that this inverse privacy loss method is the bottleneck for our implementation.

\subsection{Efficiently Representing PMoG as MoG}

As discussed in the previous section, the runtime of our implementation has a linear dependence on the number of components in the MoG. However, in \CPTB{}, we are actually using PMoGs, which are MoGs with potentially $2^n$ components. So, even just listing the components can be prohibitively expensive.

We instead choose another approximation parameter $\Delta_2$, and round each entry of $\bfC$ up to the nearest multiple of $\Delta_2$. By Lemma~\ref{lem:scalardominance}, this only worsens the privacy guarantee, i.e. any privacy guarantee we prove for the rounded version of $\bfC$ also applies to the original $\bfC$. After this rounding, the number of components in any MoG we compute the PLD of is at most $\lceil \max_i \lone{\bfe_i^\top \bfC}\rceil / \Delta_2 + n$ ($\max_i \lone{\bfe_i^\top \bfC}$ is the maximum row norm of $\bfC$). Furthermore, we can compute the probabilities/sensitivities efficiently since we are working with PMoGs. In particular, for each $\tilde{p}_{i,j}, \bfC_{i,j}$ pair, we can construct the probability mass function (PMF) of the random variable that is $\bfC_{i,j}$ w.p. $\tilde{p}_{i,j}$ and 0 otherwise, and then take the convolution of all such PMFs for a row to get the PMF of the discretized sensitivity for the PMoG. For each row, this can be done in at most $n-1$ convolutions, each convolution between two PMFs that have support size at most $2$ and $\max_i \lone{\bfe_i^\top \bfC}\rceil / \Delta_2 + n$. So the convolutions can be done in time $O(\max_i \lone{\bfe_i^\top \bfC}\rceil / \Delta_2 + n)$, i.e. our overall runtime is $O(n^2 \max_i \lone{\bfe_i^\top \bfC}\rceil / \Delta_2 + n^3)$, i.e. polynomial instead of exponential in $n$ if e.g. all entries of $\bfC$ are bounded by a constant. By doing the convolutions in a divide-and-conquer fashion, and using FFT for the convolutions, we can further improve the runtime to $\tilde{O}(n \max_i \lone{\bfe_i^\top \bfC}\rceil / \Delta_2 + n^2)$, i.e. nearly linear in the input size and $1/\Delta_2$ if the entries of $\bfC$ are bounded by a constant.

\subsection{Computing $s_{i,j}$}

Similar to computing the probabilities and sensitivities for the PMoGs, any overestimate of $s_{i,j}$ can be used in place of $s_{i,j}$ to get a valid privacy guarantee from \CPTB{} by Lemma~\ref{lem:monotonicdominance}. 
Since $s_{i,j}$ only appears in a lower order term in the definition of $\epsilon_{i,j}$, a weaker tail bound will not affect the privacy guarantee as much. So, in our implementation, we use the following simple and efficient approximation: We use the binomial CDF to obtain an exact tail bound $t$ on $\|\bfx_{1:i}\|_1 = \sum_{j' \leq i} x_{j'}$ in the definition of $s_{i,j}$. We then take the sum of the $t$ largest values of $\langle \bfC_{1:i-1,j}, \bfC_{1:i-1,j'} \rangle$ to be our overestimate of $s_{i,j}$.

\subsection{Computing All Row PLDs}
Putting this all together, we must compute $\dimension$ PLDs in \MMCC{}, one for each row of $\bfC$. Though only an $O(\dimension)$ overhead in runtime over computing a single PLD, this $O(\dimension)$ overhead is undesirable as each PLD computation is already quite expensive due to the aforementioned difficulties. However, this component is embarrassingly parallel, which we leverage to massively speed up runtimes. 

Note that for some special classes of matrices, we will have that multiple rows share the same PLD, which also allows us to dramatically speed up the calculation even without parallelization. For example, this is the case for the binary tree mechanism due to symmetry, as well for as $b$-banded Toeplitz $\bfC$ due to the fact that rows $2b-1$ to $\dimension$ of $\tilde{p}$ and $\bfC$ are the same (up to an offset in indices that doesn't affect the PLD).

\subsection{Applications Beyond Matrix Mechanisms}

We believe that MoG mechanisms/\texttt{MixtureGaussianPrivacyLoss} are useful analytic tools for privacy analysis of mechanisms beyond the matrix mechanism. We discuss two examples here.

\mypar{Privacy amplification via iteration on linear losses} Consider running DP-SGD with sampled minibatches. To get a $(\epsilon, \delta)$-DP guarantee, we can compute the PLD for the subsampled Gaussian mechanism, and then compose this PLD with itself $n$ times. For general non-convex losses, this accounting scheme is tight, even if we only release the last iterate.

For linear losses, we can give a better privacy guarantee for releasing only the last iterate, similarly to \citep{feldman2018privacy}: Releasing the last iterate is equivalent in terms of privacy guarantees to a Gaussian mechanism with random sensitivity $Binom(n, p)$ and variance $n \sigma^2$. Using \texttt{MixtureGaussianPrivacyLoss} we can get tight $(\epsilon, \delta)$-DP guarantees for this mechanism. Empirically, we found that these can be a lot tighter than composition of subsampled Gaussians. For example, using $n = 128, p = 1 / 128, \sigma = 1$ we found that composition of subsampled Gaussians gives a proof of $(.806, 10^{-6})$-DP, whereas analyzing the last iterate as a MoG mechanism gives a proof of $(.291, 10^{-6})$-DP. We conjecture a similar improvement is possible for all convex losses, rather than linear losses.

\mypar{Tight group privacy guarantees for DP-SGD} Consider analyzing the privacy guarantees of DP-SGD under group privacy. That is, we want to give a privacy guarantee for pairs of databases differing in $k > 1$ examples. One way of doing this is to compute a DP guarantee for $k = 1$, then use an example-to-group privacy theorem such as that of \cite{vadhan2017complexity}, which shows an $(\epsilon, \delta)$-DP mechanism satisfies $(k \epsilon, k \exp(k \epsilon) \delta)$-DP for groups of size $k$. This is overly pessimistic, since the black-box theorem doesn't account for the specific structure of the mechanism. We can instead get relatively tight guarantees via  \texttt{MixtureGaussianPrivacyLoss}: If each example is sampled independently, then the privacy loss of a group of $k$ examples in each round of DP-SGD is dominated by a Gaussian mechanism with sensitivity $Binom(k, p)$. The privacy loss distribution for a single instance of this mechanism can be computed using \texttt{MixtureGaussianPrivacyLoss}, and then privacy guarantees for DP-SGD follow by composition. Further, note that e.g. in the case where we instead sample a random batch of size $B$ in each round (i.e. different examples' participations within the same round are no longer independent), we can still use \texttt{MixtureGaussianPrivacyLoss} to get a tight analysis by adjusting the sensitivity random variable used. See the follow-up note \cite{ganesh2024tight} for more details. 